%%%%%%%%%%%%%%%%%%%%%%%%%%%%%%%%%%%%%%%%%%%%%%%%%%%%%%%%%%%%%%
%%%%%%%%%%%%%% MA-0794 Tópicos de Modelación %%%%%%%%%%%%%%%%%%%
%%%%%%%%%%%%%%%%%%%%%%%%%%%%%%%%%%%%%%%%%%%%%%%%%%%%%%%%%%%%%%
\newpage
\subsection*{MA-0794 Tópicos de Modelación}
\addcontentsline{toc}{subsection}{MA-0794 Tópicos de Modelación}

\hypertarget{MA0794}{}
\begin{refsection}
\begin{table}[H]
\centering{
\begin{tabular}{|ll|}
\hline
\textbf{Año:} IV  & \textbf{Requisitos:} MA-0515 Ecuaciones Diferenciales \\ 
\textbf{Ciclo:} Optativa  & \textbf{Correquisitos:} Ninguno \\ 
\textbf{Tipo de curso:} Teórico & \textbf{Horas presenciales por semana:} 4 \\ 
\textbf{Créditos:} 4 & \textbf{Horas de trabajo independiente por semana:} 8 \\ \hline
\end{tabular}
}
\end{table}

\subsubsection*{Descripción del curso}

Este curso optativo ofrece un espacio formativo para profundizar en temas de \emph{Modelación matemática} que no se desarrollan de manera sistemática en otros cursos de la carrera. Cada curso ofrecido bajo esta sigla puede orientarse hacia líneas distintas dentro de la modelación matemática---por ejemplo, pero no restringido a, aspectos clásicos o contemporáneos del área---según la especialidad de la persona docente y las necesidades formativas del estudiantado.

La naturaleza abierta del curso permite responder a avances recientes de la disciplina, a demandas del estudiantado avanzado o a líneas de investigación activas en la Escuela de Matemática, siempre dentro del marco general de la modelación matemática. Se espera que la persona estudiante consolide su capacidad de lectura de textos especializados, demostración o aplicación de resultados y comunicación matemática en un contexto de mayor autonomía intelectual.

\subsubsection*{Objetivos}

\textbf{General:} Profundizar en un tema de la modelación matemática, definido para cada oferta del curso, mediante el estudio guiado de resultados, problemas y, cuando corresponda, aplicaciones.

\textbf{Específicos:} Al finalizar el curso se espera que la persona estudiante sea capaz de:

\begin{enumerate}
\item Comprender y utilizar los conceptos centrales del bloque temático abordado en el curso correspondiente.
\item Demostrar o aplicar resultados de modelación con el nivel de rigor esperado en cursos avanzados de la carrera.
\item Consultar y sintetizar bibliografía especializada relacionada con el tema del curso.
\item Resolver problemas o desarrollar un trabajo integrador vinculado con el contenido del curso.
\item Comunicar ideas, argumentos y resultados de forma oral y escrita.
\end{enumerate}

\subsubsection*{Contenidos}

Los contenidos no se fijan de manera permanente en el plan de estudios: son \textbf{abiertos por oferta}, siempre que pertenezcan al ámbito de \emph{Modelación matemática}. Al inicio del ciclo lectivo, la persona docente debe comunicar por escrito al estudiantado:

\begin{itemize}
\item el bloque temático general que se desarrollará;
\item los objetivos específicos de esa oferta;
\item la bibliografía de referencia;
\item el cronograma de actividades y criterios de evaluación.
\end{itemize}

El curso incluirá estudio teórico, resolución de problemas y, según criterio docente, un proyecto o trabajo final integrador coherente con el tema elegido.

\subsubsection*{Metodología}

\noindent \textbf{Lineamientos metodológicos:} La persona docente a cargo de este curso debe respetar las siguientes pautas:

\begin{enumerate}
    \item Definir y publicar al inicio del curso el programa concreto de la oferta, dentro del marco general de la modelación matemática.
    \item Combinar \textbf{clases magistrales} o \textbf{seminarios} con trabajo independiente y, cuando proceda, sesiones de ejercicios o presentaciones del estudiantado.
    \item Promover la lectura activa de fuentes especializadas y la discusión crítica de resultados.
    \item Incluir evaluación formativa continua acorde con el tema y la modalidad de la oferta.
    \item Cuando las autoridades sanitarias o institucionales impongan restricciones, adaptar las actividades según normativa vigente.
\end{enumerate}

\textbf{Sugerencias metodológicas:} Adicionalmente, se tienen las siguientes recomendaciones para la persona docente a cargo de este curso:

\begin{enumerate}
    \item Coordinar con la jefatura del departamento la coherencia de la oferta con otros cursos de la modelación matemática de la malla.
    \item Explicitar los prerrequisitos conceptuales particulares de la oferta, más allá del requisito formal del curso.
    \item Fomentar la participación en coloquios, seminarios o actividades de divulgación del departamento cuando el tema lo permita.
    \item Documentar la oferta (programa, bibliografía y evaluación) para apoyar futuras reiteraciones del curso.
\end{enumerate}

\end{refsection}
